\section{Related Work}

\citet{pavel2025hallucination} characterizes hallucinations as seemingly plausible but factually incorrect or meaningless outputs that can manifest through syntax violations, invalid references, or conflicts with API knowledge, and divides them into 4 major groups: syntax, runtime, functional, and code quality hallucinations, as does the work of \citet{lee2025hallucination}. The work of \citet{dou2025whatswrong} reinforces using the first 3 classes to distinguish error types.

In our work, we distinguish the hallucination itself from its observable consequence. Instead of inferring the cause of an incorrect generation, we identify the hallucination type from its observable effect, dividing it into \emph{Syntax Error}, \emph{Execution Error}, or \emph{Functional Error}. This results-oriented abstraction is intentionally more abstract than some previous taxonomies, which, in addition to presenting macro classes, also present subclasses. However, this classification enables broader analysis of the model's deficiency rather than reducing it to a single measure of correctness. Code quality issues are outside the scope of this study, since our analysis focuses on executability, not on non-functional requirements such as style, maintainability, and efficiency.

Empirical studies have already investigated how failures vary across LLM models, including minor variations. The previously cited error-classification work, \citet{dou2025whatswrong}, also empirically evaluates, using subsets of the same benchmark employed in our article, combined across different sizes. The results showed that, regardless of size, LLMs' main difficulty is producing code that meets all functional requirements~\citep{dou2025whatswrong}. On average, smaller models were 16.03 percentage points below larger models, mainly because larger models are better at meeting these functional requirements~\citep{dou2025whatswrong}. These results provide a basis for our study, as they already describe a process for classifying error types and compare models of different sizes~\citep{dou2025whatswrong}. Our work reuses the established division and also proposes new models to understand how the distance between larger and smaller models changes over time under the same set of test cases.

The detector-oriented work most closely related to ours is \citet{andriushchenko-etal-2026-efficient}, who construct line-level hallucination annotations and train a lightweight Transformer detector using internal representations from language models. Their results demonstrate that specialized detectors can identify hallucinated regions and outperform generic uncertainty-based approaches. Our work investigates a different granularity and setting. Rather than locating individual hallucinated lines, we ask whether general-purpose and specialized SLMs can classify the observable failure mode of a complete generated solution. We additionally study self-error detection, class-specific behavior, and the interaction between model-based detection and external software-analysis evidence.

\begin{comment}
\subsection{Specialized evaluators and rationale-based supervision}

Learned evaluators have previously been incorporated into code-generation systems. CodeRL trains a critic to estimate functional correctness and uses critic predictions together with unit-test feedback to improve a code generator through reinforcement learning~\citep{le2022coderl}. Its outcome space includes compilation errors, runtime errors, failed tests, and passed tests, making it an important precedent for distinguishing different observable execution outcomes. However, the critic in CodeRL is primarily a component used to optimize generation rather than a stand-alone multilabel detector evaluated across different error categories.

Our specialized detector instead treats error recognition as the primary task. We fine-tune SLMs to predict the reference error categories directly and investigate how optimization objectives affect the resulting decision behavior. The detector is trained not only from execution-derived labels but also from natural-language explanations produced by teacher models. This design follows previous work showing that rationale-based supervision can provide richer training signals than labels alone. \citet{hsieh-etal-2023-distilling} use generated rationales as additional supervision when distilling knowledge into smaller models, while \citet{ludan-etal-2023-explanation} study explanation-based fine-tuning as a mechanism for reducing reliance on spurious cues. In HalluCodeDetection, the teacher-generated explanations serve only as auxiliary supervision: the reference labels remain determined by execution, and detector performance is evaluated exclusively from the predicted error classes.

The distinction from previous critic-based systems is therefore not the general idea of using a model to assess generated code. Rather, we investigate whether an SLM can be specialized as an independent multilabel error detector and whether specialization changes its ability to distinguish individual error categories, particularly under a strongly imbalanced class distribution.

\subsection{Self-evaluation, self-correction, and tool-assisted generation}

Several approaches improve model outputs through iterative feedback. Self-Refine alternates between generation, self-feedback, and refinement without additional model training~\citep{madaan2023selfrefine}, while Reflexion uses verbal feedback from previous attempts to guide subsequent decisions and includes programming tasks among its applications~\citep{shinn2023reflexion}. These approaches demonstrate that models can sometimes use their own assessments to revise an output, but successful revision does not necessarily imply that the underlying error was identified correctly.

For code generation, Self-Debugging asks a model to inspect execution information and reason about its program before producing a revised solution~\citep{chen2024selfdebug}. \citet{dou2025whatswrong} similarly investigate a training-free correction process based on self-critique, bug information, and compiler feedback. These studies establish that external evidence and iterative revision can improve code-generation workflows. Our study separates three related capabilities that are often considered together: recognizing an error, interpreting external evidence about that error, and successfully revising the program.

This distinction also motivates our comparison with deterministic software-analysis techniques. Compilers and static analyzers can provide direct evidence for some classes of errors, while language models may reason about behaviors that are not captured by a particular static analysis. We therefore evaluate static techniques independently, use their outputs as supporting evidence for model-based classification, and provide the same evidence during iterative code generation. This allows us to examine not only whether an error signal exists, but also whether an SLM can interpret and act on that signal appropriately.

Taken together, prior work has established taxonomies for code hallucinations, execution-based evaluation, specialized critics, self-correction, and tool-assisted generation. HalluCodeDetection builds on these directions but studies a combination that remains less explored for SLMs: generation and error detection under the same execution-based reference, self-error detection, specialization as a stand-alone multilabel detector, and the complementarity between learned detection and deterministic software-analysis evidence. This design allows us to investigate whether generation, error recognition, evidence interpretation, and correction behave as the same capability or as distinct aspects of reliable code generation.
\end{comment}